\documentclass[11pt,a4paper]{article}

%----- ENHANCED TYPOGRAPHY -----
\usepackage[utf8]{inputenc}
\usepackage[T1]{fontenc}
\usepackage{lmodern}        % clean vector font
\usepackage{microtype}      % better justification & kerning
\usepackage{palatino} 
\usepackage{braket}   
\usepackage{mathtools}  % in the preamble
\usepackage{bbm}

%----- PAGE LAYOUT -----
\usepackage{geometry}
\geometry{top=1in, bottom=1in, left=1in, right=1in}
\usepackage{setspace}
\onehalfspacing  % 1.5 line spacing

%----- FANCY HEADERS & FOOTERS -----
\usepackage{fancyhdr}
\pagestyle{fancy}
\fancyhf{}
% page number outside, header text inside

\fancyhead[LO]{\small \rightmark}
\fancyhead[RO]{\small \leftmark}
\renewcommand{\headrulewidth}{0.4pt}
\renewcommand{\footrulewidth}{0pt}

% make sections feed into \leftmark/\rightmark
\renewcommand{\sectionmark}[1]{\markboth{#1}{}}
\renewcommand{\subsectionmark}[1]{\markright{#1}}

%----- SECTION NUMBERING & TOC DEPTH -----
\setcounter{secnumdepth}{3}  % number down to \subsubsection
\setcounter{tocdepth}{2}     % show ToC down to \subsection

%----- AMS MATH & THEOREM STYLES -----
\usepackage{amsmath,amssymb,mathtools}
\usepackage{amsthm}

% definitions, examples, remarks upright
\theoremstyle{definition}
\newtheorem{definition}{Definition}[section]
\newtheorem{example}[definition]{Example}
\newtheorem{remark}[definition]{Remark}

% theorems, lemmas, corollaries italic
\theoremstyle{plain}
\newtheorem{theorem}[definition]{Theorem}
\newtheorem{lemma}[definition]{Lemma}
\newtheorem{proposition}[definition]{Proposition}
\newtheorem{corollary}[definition]{Corollary}

% unnumbered proof environment
\theoremstyle{remark}

%----- OTHER PACKAGES -----
\usepackage{graphicx}
\usepackage{tikz}
\usetikzlibrary{arrows.meta,positioning}
\usetikzlibrary{calc, matrix, decorations.pathreplacing, positioning}
\usepackage{tikz-cd}
\usepackage{hyperref}
\hypersetup{colorlinks,
linkcolor=blue, citecolor=purple, urlcolor=teal}
\usepackage{enumitem}
\setlist[itemize]{nosep, left=1.5em}
\usepackage{booktabs}
\usepackage{listings}
\lstset{
basicstyle=\ttfamily\small,
numbers=left,
numbersep=5pt,
frame=single,
breaklines=true
}
\usepackage{xcolor}
\definecolor{shade}{HTML}{F5F5F5}
\usepackage{float}
%----- CUSTOM MACROS -----
\newcommand{\F}{\mathbb{F}}
\newcommand{\code}[1]{\texttt{#1}}
\newcommand{\dist}[2]{d\bigl(#1,#2\bigr)}
\newcommand{\R}{\mathbb{R}}
\newcommand{\Z}{\mathbb{Z}}
\newcommand{\N}{\mathbb{N}} 
\newcommand{\Q}{\mathbb{Q}} 
\newcommand{\C}{\mathbb{C}}
\renewcommand{\set}[1]{\left\{ #1 \right\}}
\newcommand{\angles}[1]{\langle #1 \rangle}
\newcommand{\abs}[1]{\lvert #1 \rvert}
\newcommand{\norm}[1]{\lVert #1 \rVert}
\newcommand{\1}{\mathbbm{1}}

% \usepackage{mathtools} 
% \DeclarePairedDelimiter{\angles}{\langle}{\rangle} 
% \DeclarePairedDelimiter{\braces}{\left\{}{\right\}} 
% \DeclarePairedDelimiter{\abs}{\lvert}{\rvert} 
% \DeclarePairedDelimiter{\norm}{\lVert}{\rVert}

%----- TITLE METADATA -----
\title{\LARGE\bfseries Algorithmic and High Frequency Trading}
\author{Georges Khater}
\date{\today}

%===============================================
\begin{document}
\maketitle
\tableofcontents
\bigskip

\section{Electronic Markets and the Limit Order Book.}

\subsection{Electronic Markets and How they Function.} 
Review basic financial definitions. 

\subsection{Classifying Market Participants.}
When designing Algorithmic trading, it is important to consider who trades in the exchanges and why. Since 
this will affect whether and how algorithms achieve the desired trading objective.  
\begin{itemize}
    \item Fundamental traders (noise / liquidity traders). 
    \item Informed traders (active or aggressive). 
    \item Market Makers (passive or reactive trading). 
\end{itemize}

\subsection{Trading in Electronic Markets} 
Electronic markets amount to having a way for people to show their willingness to trade with an engine to match those 
wanting to buy with those wanting to sell. 

\subsubsection{Orders and the Exchange}
There are two types of orders
\begin{itemize}
    \item Market orders; wants to buy or sell a quantity of shares 
    at the best available price. Often results in immediate trade execution.

    \item Limit orders; passive trades orders which indicate willingness to 
    buy or sell at a given price, up to a maximum quantity of shares. 
    This usually doesn't result in an immediate trade; it will have to wait until it is matched with some other 
    order or is withdrawn. 
\end{itemize}
These trades are managed by a matching engine and a limit order book (LOB). 
The latter keeps track of incoming and going orders. Matching engine is a well defined 
algorithm to match sellers to buyers. 

[insert picture] 
Notice how the difference in liquidity affects the LOBs. 

\subsubsection{Alernate Exchange Structure.}
It is also important to take into consideration the degree (and cost) of transparency in 
an exchange. We distinguish lit (open OB) from dark markets. 

\subsubsection{Colocation.} 
Exchanges control the amount of information you receive, they also rent out 
computer space next to their matching engines. This process is called colocation. 

We also have to keep in mind the interaction of multiple trading venues, 
latency when moving between these venues and regulation when designing trading strategies. 

\subsubsection{Extended Order Types.} 
The race to be the first in or out of a certain 
position is one of the focal points of the debate on the benefits and costs of 
HFT.

This importance of speed affects the whole process of designing trading algorithms 
(code, language, hardware, etc).  Hence exchanges have moved well-beyond the basic types of 
orders (MOs and LOs): 
\begin{itemize}
    \item Day orders: orders for trading during regular trading with options to extend
    to pre- or post-market sessions. 

    \item Non-routable: there are a number of orders that by choice or design avoid
    the default re-routing to other exchanges, such as 'book only', 'post only', 
    'midpoint peg', etc. 

    \item Pegged, Hide-not-Slide: orders that move with the midpoint or the national
    best price

    \item Hidden: orders that do not display their quantity
    \item Iceberg: orders that partially display their quantity (some have options so that the visible portion will automatically be replenished when it is depleted 
    by less than one round lot)
    
    \item Immediate-or-Cancel: orders that execute as much as possible at the best
    price and the rest are cancelled (such orders are not re-routed to another exchange nor do they walk the book)
    
    \item Fill-or-Kill: orders sent to be executed at the best price in their entirety or
    not at all; 

    \item Good-Till-Time: orders with a fixed lifetime built into them so that they
    will be cancelled if not executed by its expiration time
    
    \item Discretionary: orders display one price (the limit price) but may be executed at more aggressive (hidden) prices; 
\end{itemize}
Note that we should be aware of the different types of orders not only in our exchange, but across different exchanges. 

\subsubsection{Exchange Fee}
Trading in an exchange is not free; usually a much higher fee on MOs, sometimes LOs are due rebates, or the reverse in some markets.
This distorts observed market prices. 

\subsection{The Limit Order Book.} 
$$\text{Quoted Spread}_t = P_t^a - P_t^b$$
if the spread is zero the market becomes locked. 

$$\text{Midprice}_t = \frac{1}{2} (P_t^a + P_t^b)$$
usually used to proxy the price of teh asset. 

[insert picture]
Note the difference in the LOB of liquid vs illiquid asset. 

When a sell MO executes against a buy LO, it is said to \emph{hit the bid}; 
otoh when a buy MO executes against a sell LO it is said to \emph{lift the offer}. 

$$\operatorname{Microprice}_t = \frac{V_t^b}{V_t^b + V_t^a} P_t^a + \frac{V_t^a}{V_t^b + V_t^a} P_t^b$$
Where $V_t^b$ and $V_t^a$ are the volumes posted at the best bid and asks. 

\subsection*{A Primer on the Microstructure of Financial Markets.} 
The goal is to consider the economics that drive these trading strategies. 

For the economics of trading we look to market microstructure, as it is the 
subfield of finance which focuses on how trading takes place in very specific 
setting. 

A key dimension of the trading and price setting process is that of information. Who has what information, how does that information affect trading strategies, 
and how do those trading strategies affect trading outcomes in general, and asset prices in particular.

What differentiates microstructure studies from more general 
asset pricing ones is that they focus on two aspects that are key to trading: 
liquidity and price discovery.

\subsection*{Market Making.} 
The market makers provide liquidity in exchange for spread, we will consider that there are a lot of market makers. 

\subsubsection{Grossman-Miller Market Making Model}
The main risk of the above is that when accepting one side of a trade, the MM will be exposed for an undetermined period of time. 
Grossman \& Miller (1988) provide a model that captures this problem and describes how 
MMs obtain a liquidity premium from liquidity traders that exactly compensates 
MMs for the price risk of holding an inventory of the asset until they can unload 
it later to another liquidity trader. 

Read the section in the book describing the model, nice stuff. 

The size of the liquidity shock, risk aversion, 
and volatility all increase the discount, while competition reduces it. 

\subsubsection{Trading Costs.} 
We saw above how the costs of holding assets affect liquidity via the costs of trading. But what drives $n$ (the number of MMs)? 
We do this by linking $n$ with participation costs. This decision to participate is determined as a function of a participation cost parameter
$c$. 

We are also interested in the effect of a cost of trading that depends on the level of activity in the market. 
We introduce trading costs that depend on the quantity traded, like actual exchange trading fees. By introducing such a 
parameter $\eta$ assuming that all $|i|$ is big enough for trades to be desirable we get that the liquidity discount is bigger. 

Liquidity traders incur most of the trading fees ($(1 + 2 \frac{n}{n+1})\eta$), and receives less immediacy from the market. 

Therefore an increase in trading has a smaller effect on liquidity via competition but a larger effect on immediacy and liquidity discount. 

\subsubsection{Measuring Liquidity}
We now consider how these divergences from efficient prices may be observed in electronic exchanges. 
Note that we should (and it is easy to) adapt the analysis above to the case of an electronic market. 

The Grossman-Miller model corresponds to the following sequence: as LT1's MOs enter the market, they execute against the MM's LOs in the 
LOB which adjusts to the incoming MO. As the execution price changes, so does the LT1's strategy who eventually stops trading after having sold 
$i \frac{n}{n+1}$ shares. $S_1$ in this case is the average price of the executed trades. 

We can write $S_1$ as a linear function of the quantity traded $q^{LT1}$: 
$$S_1 = \mu + \lambda q^{LT1}$$
so that in the Grossman-Miller model we would have 
$$\lambda = -\frac{1}{n} \gamma \sigma^2, \quad q^{LT1} = i \frac{n}{n+1}$$
$\lambda$ here captures the market's price impact. In particular, $\lambda$ is used as a proxy for liquidity of the market. 
A more liquid market will have a lower $\lambda$, either because of greater competition ($n$), lower risk tolerance $\gamma$ or lower 
volatility $\sigma^2$. 

Another way to measure liquidity based on price changes is the autocovariance of the asset's returns. 
Let $\Delta_1 = S_1 - S_0$ and $\Delta_2 = S_2 - S_1$, then the autocovariance of price changes is given by the following relation: 
\begin{align*}
    \operatorname{cov} [\Delta_1, \Delta_2] &= \operatorname{Cov} [\mu_1 + \lambda q^{LT1} - \mu_0, \mu_2 - \mu_1 - \lambda q^{LT1}] \\
    &= - \lambda^2 \operatorname{Var} [q^{LT1}] < 0
\end{align*}

\subsubsection{Market Making using Limit Orders.}
We proposed above that MMs would participate through posting LOs, we know explain why this is a natural assumption to make. 
Consider a small, risk-neutral trader with costless inventory and infinite patience; she makes profit by optimally providing liquidity 
in an uncertain environment populated by other MMs who do not react to our MM's decisions. 

Uncertainty in this context comes from the timing and size of large incoming 
MOs, and there are no information problems: all information is public so that 
everyone agrees what the current value of the asset is, which we denote by St 
and refer to as the midprice.

We wat to choose the distance from the midprice (the depths $\delta^{\pm}$) that maximize profit. The probability 
that once an MO arrives it will walk up the book to a specific price is described by $P_{\pm}$. 
i.e. the probability that the buy LO will be filled is $p_- P_- (\delta^-)$. We assume that the 
distribution of other LOs in the LOB is described by an exponential distribution with parameter $\kappa^-$. 
This gives us that 
$$p_-P_-(\delta-) = p_- e^{- \kappa^- \delta^-}$$

Using $\Pi$ to denote the MM's profit per trade, we get that the MM's optimization problem is given by 
$$\max_{\delta^+, \delta^-} \mathbb{E} [\Pi(\delta^+, \delta^-)] = \max_{\delta^+, \delta^-} \set{p^+ e^{- \kappa^+ \delta^+} \delta^+ + p^- e^{- \kappa^- \delta^-} \delta^-}$$
Which gives the solutions 
$$\delta^{\pm} = \frac{1}{\kappa^\pm}$$
Note that this model is very unrealistic, the functional form of all the stochastic components of the model ($P_\pm$ and $p^\pm$)  
are constant and exogenous. The MM's decisions and that are other trade ($P(\delta)$) are independent, the MM's objective function is static and very simple. 

\subsection*{Trading on an Informational Advantage}

\end{document}